%==============================================================================
%  MartFlow --- Supershop Retail Management Suite
%  CSE 464 : Software Design Pattern Lab  |  Summer 2026
%  Final Project Report
%
%  Compile on Overleaf with pdflatex; run 2-3 passes for TOC / LOF / LOT.
%
%  DIAGRAM WORKFLOW (Mermaid):
%   - Every figure below that uses \dfig has its Mermaid source in a comment
%     block directly above it, opened by a "MERMAID: <slug>" banner.
%   - Render at https://mermaid.live (strip the leading "% "), export PNG 3x,
%     save as diagrams/<slug>.png, upload into a "diagrams" folder in
%     Overleaf. The placeholder box is replaced by the real image
%     automatically -- no .tex edits needed.
%   - UI screenshots go into a "figs" folder (see the \uifig slots in
%     Chapter 5); missing files render a placeholder and never break the
%     build.
%==============================================================================
\documentclass[12pt,a4paper]{report}

%---------- Geometry & page setup ----------
\usepackage[a4paper,top=1.0in,bottom=1.0in,left=1.1in,right=1.1in,headheight=15pt]{geometry}
\usepackage{fancyhdr}
\usepackage{titlesec}
\titleformat{\chapter}[display]
  {\normalfont\huge\bfseries}{\chaptertitlename\ \thechapter}{10pt}{\Huge}
\titlespacing*{\chapter}{0pt}{-60pt}{20pt}

%---------- Encoding / fonts (pdflatex-safe, 7-bit ASCII diagrams) ----------
\usepackage[T1]{fontenc}
\usepackage[utf8]{inputenc}
\usepackage{lmodern}

%---------- Math, tables, graphics ----------
\usepackage{amsmath,amssymb}
\usepackage{graphicx}
\usepackage{booktabs}
\usepackage{array}
\usepackage{tabularx}
\usepackage{longtable}
\usepackage{multirow}
\usepackage{float}
\usepackage{caption}
\captionsetup{font=small,labelfont=bf}

%---------- Colour, lists, code, links ----------
\usepackage{xcolor}
\usepackage{enumitem}
\setlist[itemize]{leftmargin=*,topsep=2pt,itemsep=1pt}
\setlist[enumerate]{leftmargin=*,topsep=2pt,itemsep=1pt}
\usepackage{listings}
\usepackage{url}
\usepackage[hidelinks]{hyperref}

%---------- TikZ for diagrams ----------
\usepackage{tikz}
\usetikzlibrary{positioning,arrows.meta,shapes.geometric,shapes.misc,calc,fit,backgrounds}

%---------- Brand palette (emerald + gold, matching the app) ----------
\definecolor{emerald}{HTML}{047857}
\definecolor{emeraldlight}{HTML}{10B981}
\definecolor{gold}{HTML}{C99A2E}
\definecolor{slate}{HTML}{475569}
\definecolor{codebg}{HTML}{F6F8F7}

%---------- Listings styles (code + ASCII diagrams) ----------
\lstdefinelanguage{JavaScript}{
  morekeywords={const,let,var,function,return,if,else,for,await,async,new,try,catch,require,module,exports,true,false,null,typeof},
  morecomment=[l]{//},
  morestring=[b]',
  morestring=[b]"
}

\lstdefinelanguage{Java}{
  morekeywords={abstract,assert,boolean,break,byte,case,catch,char,class,continue,default,do,double,else,enum,extends,final,finally,float,for,if,implements,import,instanceof,int,interface,long,native,new,package,private,protected,public,return,short,static,super,switch,synchronized,this,throw,throws,transient,try,void,volatile,while,true,false,null,var,record},
  sensitive=true,
  morecomment=[l]{//},
  morecomment=[s]{/*}{*/},
  morestring=[b]",
  morestring=[b]'
}

\lstset{
  basicstyle=\ttfamily\footnotesize,
  backgroundcolor=\color{codebg},
  breaklines=true,
  frame=single,
  framerule=0.3pt,
  rulecolor=\color{slate!40},
  keywordstyle=\color{emerald}\bfseries,
  commentstyle=\color{slate}\itshape,
  stringstyle=\color{gold!80!black},
  showstringspaces=false,
  numbersep=6pt,
  xleftmargin=10pt
}

\lstdefinestyle{ascii}{
  language=,
  basicstyle=\ttfamily\scriptsize,
  frame=single,
  numbers=none,
  keywordstyle=,
  commentstyle=,
  stringstyle=,
  backgroundcolor=\color{codebg}
}

%---------- Headers / footers ----------
\pagestyle{fancy}
\fancyhf{}
\fancyhead[L]{\small\itshape MartFlow}
\fancyhead[R]{\small Software Design Pattern Lab}
\fancyfoot[C]{\thepage}
\renewcommand{\headrulewidth}{0.4pt}

\fancypagestyle{plain}{%
  \fancyhf{}
  \fancyfoot[C]{\thepage}
  \renewcommand{\headrulewidth}{0pt}
}

%---------- Reference / appendix names ----------
\renewcommand{\bibname}{References}
\renewcommand{\contentsname}{Table of Contents}

%---------- Helper: Mermaid diagram figure with graceful placeholder ----------
% \dfig{Caption}{fig:label}{diagrams/slug.png}{width-fraction}
% The Mermaid source sits in a comment block IMMEDIATELY above each call.
\newcommand{\dfig}[4]{%
  \begin{figure}[H]
  \centering
  \IfFileExists{#3}{%
    \includegraphics[width=#4\textwidth,keepaspectratio]{#3}%
  }{%
    \setlength{\fboxsep}{10pt}\setlength{\fboxrule}{0.4pt}%
    \fcolorbox{slate!40}{codebg}{%
      \begin{minipage}[c][6.2cm][c]{0.88\textwidth}\centering
        \itshape\small Render the Mermaid block above this figure at\\
        \url{https://mermaid.live} and export a 3$\times$ PNG as\\[2pt]
        \texttt{\detokenize{#3}}\\[2pt]
        This box is then replaced by the real diagram automatically.
      \end{minipage}}%
  }%
  \caption{#1}\label{#2}
  \end{figure}}

%---------- Helper: UI screenshot loader / placeholder ----------
% \uifig{Caption}{fig:label}{figs/name.png}
\newcommand{\uifig}[3]{%
  \begin{figure}[H]
  \centering
  \IfFileExists{#3}{%
    \includegraphics[width=0.88\textwidth]{#3}%
  }{%
    \setlength{\fboxsep}{10pt}
    \setlength{\fboxrule}{0.4pt}%
    \fcolorbox{slate!40}{codebg}{%
      \begin{minipage}[c][5.2cm][c]{0.88\textwidth}
        \centering
        \itshape\small \detokenize{#3}
      \end{minipage}
    }%
  }%
  \caption{#1}\label{#2}
  \end{figure}}

\begin{document}

%==============================================================================
%  COVER PAGE
%==============================================================================
\begin{titlepage}
\thispagestyle{empty}

\vspace*{0.4cm}

\noindent
\begin{minipage}{0.15\textwidth}
    \IfFileExists{bubtlogo.png}{%
        \includegraphics[width=\linewidth]{bubtlogo.png}%
    }{%
        \framebox[\linewidth]{%
            \rule{0pt}{0.8\linewidth}
            \centering\small BUBT
        }%
    }
\end{minipage}%
\begin{minipage}{0.85\textwidth}
    \centering
    \large \textbf{Bangladesh University of Business and Technology}\\[1ex]
    \large Department of Computer Science and Engineering
\end{minipage}

\vspace{1.5cm}

\begin{center}
    \Large CSE 464 : Software Design Pattern Lab
\end{center}

\vspace{1cm}

\begin{center}
    \Huge \textbf{PROJECT REPORT}
\end{center}

\vspace{1.5cm}

\noindent
\begin{tabular}{@{}l @{\hspace{0.4em}:\hspace{0.8em}} l@{}}
    \textbf{Semester}        & Summer 2026 \\[1.2ex]
    \textbf{Title}           & MartFlow --- Supershop Retail Management Suite \\[1.2ex]
    \textbf{Submission Date} & 17 August 2026
\end{tabular}

\vspace{1.5cm}

\begin{center}
    \textbf{TEAM MEMBERS}\\[2ex]

    \renewcommand{\arraystretch}{1.3}

    \begin{tabular}{|c|c|c|}
        \hline
        \textbf{Full ID} & \textbf{Full Name} & \textbf{Intake/Section} \\
        \hline
        22235103137 & Md. Mahamudul Hasan & 51/11 \\
        \hline
        22235103251 & Sumya Soma & 51/11 \\
        \hline
        22235103680 & Alfa Sany & 51/11 \\
        \hline
    \end{tabular}
\end{center}

\vspace{2cm}

\begin{center}
    \textbf{SUPERVISED BY}\\[2ex]

    \textbf{Talimul Bari Shreshtho}\\
    Lecturer, Department of CSE\\
    Bangladesh University of Business and Technology (BUBT)
\end{center}

\end{titlepage}

%==============================================================================
%  FRONT MATTER
%==============================================================================
\pagenumbering{roman}

%---------- Abstract ----------
\chapter*{Abstract}
\addcontentsline{toc}{chapter}{Abstract}
\noindent
MartFlow is a supershop retail management suite that gives a Bangladeshi supershop branch a
single system for point-of-sale billing, inventory, purchasing, promotions, loyalty, returns
and financial reporting. The presentation tier is a vanilla-JavaScript single-page
application served by the back end: a hash router drives 15 role-gated screens styled with
Bootstrap~5.3.3, including a thermal-receipt printer. The application tier is written in
Java~17 on Spring Boot~3.4.0 and deliberately keeps the domain 100\%~framework-free: 15 REST
controllers expose 62 endpoints, each delegating to plain-Java facades and services in which
18 GoF design patterns (4~creational, 5~structural and 9~behavioral) each carry a real
business feature. The data tier is MongoDB Atlas with 11~collections, behind a generic
repository interface with a zero-configuration in-memory fallback for demos. Authentication
uses PBKDF2 password hashing and 256-bit bearer tokens with a sliding 12-hour lifetime;
authorization is enforced in three layers --- hidden menus, server-side role gates on every
controller, and a role-guard proxy on the data boundary --- across four roles: owner,
manager, cashier and developer. The platform delivers 16 functional modules, including
VAT-inclusive billing with NBR slab back-calculation (0\%, 7.5\%, 15\%), split tenders over
cash, card, bKash, Nagad and loyalty points, a guarded purchase-order state machine, partial
returns with pro-rata refunds, a reconciled day-close Z-report, an append-only audit trail,
and eight report types exportable to CSV. Every money figure is an exact-decimal
\texttt{BigDecimal}, all business time is pinned to Asia/Dhaka, and 176 automated tests in
31 hermetic test classes protect the money paths. A built-in Developer Mode documents each
pattern with its real source snippet and proving test, making the suite a working product
and a living design-patterns laboratory at the same time.

\vspace{1em}
\noindent\textbf{Keywords:} point of sale, supershop retail, Bangladesh, NBR VAT,
inventory, purchasing, loyalty, GoF design patterns, Java, Spring Boot, MongoDB Atlas,
PBKDF2, role-based access control, single-page application.

%---------- Acknowledgment ----------
\chapter*{Acknowledgment}
\addcontentsline{toc}{chapter}{Acknowledgment}
\noindent
We express our sincere gratitude to our supervisor, \textbf{Talimul Bari Shreshtho},
Lecturer, Department of Computer Science and Engineering, Bangladesh University of Business
and Technology, for continuous guidance, valuable suggestions and encouragement throughout
the development of this project. We are also thankful to the faculty members of the
Department for their support, and to the open-source community --- the Java, Spring Boot,
MongoDB, Bootstrap and JUnit projects in particular --- whose libraries and documentation
made this work possible. Finally, we thank our families for their patience and
encouragement during the entire CSE~464 cycle.

%---------- Table of contents / list of figures / list of tables ----------
\tableofcontents
\cleardoublepage
\addcontentsline{toc}{chapter}{List of Figures}
\listoffigures
\cleardoublepage
\addcontentsline{toc}{chapter}{List of Tables}
\listoftables
\cleardoublepage

%==============================================================================
%  MAIN MATTER
%==============================================================================
\pagenumbering{arabic}

%==============================================================================
\chapter{Introduction}
%==============================================================================

\section{Problem Specification}
Bangladeshi supershops --- the Shwapno, Meena Bazar and Unimart class, plus hundreds of
mid-size local chains --- typically run their branches on a patchwork of tools. Imported
point-of-sale software does not understand Bangladesh VAT: shelves here display
VAT-inclusive prices, output VAT must be back-calculated per National Board of Revenue
(NBR) rate slab (0\%, 7.5\%, 15\%), and foreign products price in dollars per month.
Purchasing lives in spreadsheets, loyalty lives in notebooks, and wastage is tracked
nowhere at all. The owner finds out about problems weeks late: stock that walked out the
door, batches that expired on the shelf, and margins that quietly eroded. Local cash
registers record sales but cannot answer ``who voided this bill?'', cannot reconcile the
drawer at day close, and cannot produce the VAT summary the tax office expects.

The problem is sharpened by what a course project usually produces. An earlier e-commerce
demonstration built by this team exhibited the classic flaws: no real login, the role taken
from a URL parameter, one shared cart for every visitor, a tax figure that was displayed
but never charged, and stock that evaporated on every restart. Each of those flaws
represents a class of real business failure --- unauthorized access, privilege escalation,
data races, phantom revenue and data loss.

These gaps motivated MartFlow, a single suite in which the money path is honest by
construction: prices are VAT-inclusive with per-slab back-calculation, every amount is an
exact decimal, every void and return is fully reversed and audited, and the drawer
reconciles by hand at day close.

\section{Objectives}
The principal objectives of the project were:
\begin{enumerate}
  \item To build a complete supershop management system with a vanilla-JavaScript
        single-page front end, a Java~17 / Spring Boot~3.4.0 application layer and a
        MongoDB Atlas data layer with an in-memory fallback.
  \item To implement point-of-sale billing that matches how Bangladeshi shelves actually
        work: VAT-inclusive pricing with NBR slab back-calculation, per-piece and weighed
        goods, combo hampers, and split tenders over cash, card, bKash, Nagad and loyalty
        points.
  \item To keep inventory truthful: cost and MRP per item, batch and expiry tracking,
        reorder levels, and damage/loss/theft write-offs (shrinkage) as first-class events.
  \item To manage the purchasing cycle end to end: suppliers with payment terms, purchase
        orders with a guarded state machine, goods receipts that land stock with batch and
        updated cost, payables, and weekly restock templates.
  \item To enforce exact-decimal money (\texttt{BigDecimal}, scale~2, \texttt{HALF\_UP})
        through a single choke point, and business time pinned to Asia/Dhaka.
  \item To secure the platform with PBKDF2-hashed passwords, bearer tokens and four roles
        (owner, manager, cashier, developer) enforced in three layers, from the menu down
        to the data boundary.
  \item To make the shop accountable: an append-only activity log of every sensitive
        action, and a day-close Z-report whose expected-drawer formula reconciles by hand.
  \item To apply 18 GoF design patterns --- and only where a real feature needs each ---
        protected by 176 hermetic automated tests, and to expose that reasoning inside the
        product through a dedicated Developer Mode.
\end{enumerate}

\section{Scope}
The delivered system covers staff authentication and role-based access control; POS
billing with undo, coupons, charges and split tender; inventory with batches, expiry,
reorder views and shrinkage adjustments; the purchasing cycle with suppliers, purchase
orders, goods receipts, payments, closing and standing restock templates; promotions
(category sales, member pricing and coupon codes) with a live tester; a loyalty program
with points earned and redeemed at the till; partial per-line returns with pro-rata
refunds and manager-side voids; a sales explorer with reprint; a day-close Z-report with
drawer variance; automatic alerts with a suggestion service; an append-only activity log;
eight report types with CSV export and an operational dashboard; staff management; and a
Developer Mode for pattern inspection, API exploration and diagnostics.

Out of scope for this version, and discussed under future work, are: native mobile
applications; a multi-branch synchronisation server; real merchant integrations with bKash,
Nagad or card acquirers (vendor-shaped SDKs are simulated behind adapters); offline-first
billing with background reconciliation; barcode-printer and cash-drawer hardware; and
machine-learning analytics. Client-side internet access is required for the Bootstrap CDN.

\section{Flow of the Project}
Development was organized into six phases:
\begin{description}[leftmargin=2.2em]
  \item[Phase 1 --- Planning \& Design (Weeks 1--2):] Requirements gathered from real
        supershop workflows (billing, purchasing, day close), the 18-pattern mapping was
        drawn, the entity-relationship model of the 11 collections was designed, and UI
        wireframes were sketched.
  \item[Phase 2 --- Domain \& Back-End Development (Weeks 3--5):] The framework-free
        domain (billing, catalog, inventory, pricing, payment, suppliers, reports) with
        its design patterns, the persistence layer with Atlas and in-memory repositories,
        PBKDF2 authentication, and all 15 REST controllers.
  \item[Phase 3 --- Front-End Development (Weeks 6--9):] The SPA shell, hash router with
        its role gate, the POS screen, and the management screens (inventory, purchasing,
        loyalty, returns, reports, sales explorer).
  \item[Phase 4 --- Feature Enhancement (Weeks 10--11):] The accountability layer:
        day-close Z-report, audit trail, promotions manager, sales explorer and Developer
        Mode.
  \item[Phase 5 --- Testing \& QA (Weeks 12--13):] The hermetic 176-test suite, restart
        drills against Atlas (byte-identical receipt reprints), drawer-math checks and
        security hardening.
  \item[Phase 6 --- Documentation \& Presentation (Weeks 14--15):] README and Bengali
        companion documentation, drift guards, this report and the presentation.
\end{description}

\section{Organization of the Project Report}
The remainder of this report is organized into six chapters.
\textbf{Chapter~2} surveys existing systems and the supporting literature.
\textbf{Chapter~3} presents the system analysis and design --- the technology stack, the
development model, the system architecture, the applied GoF design patterns, the use-case,
context and data-flow diagrams, the database schema, and the core algorithms.
\textbf{Chapter~4} describes the implementation of the front end and back end and
enumerates the functional modules. \textbf{Chapter~5} is a user manual covering system
requirements and interface walkthroughs. \textbf{Chapter~6} concludes with a discussion of
limitations and future work, followed by the references and the appendices
(Appendix~\ref{app:glossary}: glossary, key API endpoints and the project timeline).

%==============================================================================
\chapter{Background and Related Works}
%==============================================================================

\section{Existing System Analysis}
Several classes of software already serve parts of the supershop problem, but each leaves
gaps that a Bangladeshi branch feels daily. Table~\ref{tab:competitors} summarises the
most relevant of them.

\begin{table}[H]
\centering
\caption{Comparative analysis of existing platforms.}
\label{tab:competitors}
\small
\renewcommand{\arraystretch}{1.25}
\begin{tabularx}{\textwidth}{l X X}
\toprule
\textbf{Platform} & \textbf{Strengths} & \textbf{Weaknesses} \\
\midrule
Shopify POS & Polished till UX, huge app ecosystem, omnichannel inventory &
Dollar-priced subscription; no NBR VAT slabs or VAT-inclusive back-calculation; no bKash /
Nagad / points tenders \\
Square POS & Clean hardware and offline-till pedigree, strong analytics &
Not available in Bangladesh; foreign tax model; no local payment rails \\
Lightspeed Retail & Deep inventory with purchasing and vendor management &
Expensive per-month pricing; cloud-only; BD VAT and drawer-reconciliation workflows absent \\
TallyPrime (ERP class) & Accounting-grade ledgers, used widely in Bangladesh,
tax summaries & Till-first workflow weak; clerk training burden; POS, loyalty and returns
feel bolted on \\
Local legacy POS software & Cheap, offline, Bengali-friendly &
No batch/expiry or shrinkage tracking; no day-close reconciliation; little or no audit
trail; stale UIs; purchasing still in spreadsheets \\
\bottomrule
\end{tabularx}
\end{table}

The recurring gaps identified were:
\begin{itemize}
  \item No VAT-inclusive pricing with NBR slab back-calculation (0\%, 7.5\%, 15\%) feeding
        a filing-ready VAT summary.
  \item No split tender across cash, card, bKash, Nagad and loyalty points in one bill.
  \item No single system joining purchasing (with a guarded order lifecycle) to the shelf,
        the till, returns, and a reconciled day close.
  \item No accountability layer: an append-only audit trail with the actor behind every
        void, return and price edit, and a Z-report whose drawer math reconciles by hand.
  \item No affordable, zero-configuration option a mid-size branch can actually run ---
        and no transparency about the engineering underneath.
\end{itemize}

MartFlow addresses every one of these gaps in one suite, and its Developer Mode makes the
design-pattern reasoning behind each decision inspectable from inside the running product.

\section{Supporting Literatures}
The work draws on several bodies of knowledge --- Bangladeshi retail and VAT practice, the
classic design-patterns catalogue, layered web architecture, authentication frameworks, and
the specific front-end and back-end technologies employed.

\paragraph{Retail domain and Bangladesh VAT.} The Value Added and Supplementary Duty Act
2012 governs the slabs applied here; supershop shelves display VAT-inclusive prices, so the
system back-calculates output VAT per line using
$\mathrm{VAT} = \mathrm{net} \times r/(100+r)$ in one shared \texttt{VatCalculator}, and
the VAT summary report groups taxable value and output VAT by slab --- the format a filing
expects. The day-close Z-report mirrors the physical practice of drawer reconciliation at
shift end.

\paragraph{Design patterns.} The Gang-of-Four catalogue~\cite{gof} and its modern
teaching treatments~\cite{headfirst,refactoringguru} supplied the vocabulary. The project's
rule is stricter than the syllabus: a pattern earns its place only by carrying a real
feature. Three patterns that were dead code in the earlier demonstration (Prototype, a tax
visitor, a customer observer) were rebuilt into load-bearing features --- weekly restock
templates, the NBR VAT report and live loyalty points respectively --- and several
candidates were deliberately rejected (Section~\ref{sec:patterns}).

\paragraph{Architecture.} The system follows a classical three-tier pattern with a strict
layer rule: a browser-served SPA presentation layer, an application layer of thin Spring
controllers over plain-Java facades and services, and a data layer behind a generic
repository interface~\cite{cleanarch,sommerville}. The domain contains no framework
annotations at all; Spring appears only in the controllers, one filter and one
configuration class that assembles the object graph --- which is exactly what keeps the
patterns honest and the domain unit-testable.

\paragraph{Authentication and authorization.} Passwords are hashed with PBKDF2 as
specified in RFC~8018~\cite{rfc8018} (120{,}000 iterations, per-user salt, constant-time
comparison), following the OWASP password-storage guidance~\cite{owasp}. Sessions are
256-bit random bearer tokens with a sliding 12-hour lifetime. Authorization follows
role-based access control enforced server-side; the caller's role never travels in the
request body, and a ThreadLocal request context is always cleared in a \texttt{finally}
block~\cite{owasp,fielding}.

\paragraph{Front-end stack.} The user interface is deliberately framework-free: vanilla
JavaScript with a hash router (no server route coupling, trivial hosting), the Fetch API
with a Bearer header~\cite{mdnjs,mdnfetch}, and Bootstrap~5.3.3~\cite{bootstrap} with
Bootstrap Icons~\cite{bootstrapicons} for a familiar, responsive component vocabulary.

\paragraph{Back-end stack and persistence.} Spring Boot~3.4.0~\cite{springboot} supplies
the embedded server, serialization and test tooling; the MongoDB Java
driver~\cite{mongodriver} talks to MongoDB Atlas~\cite{mongodb,atlas} through hand-written
repositories and document mappers rather than a heavy data framework, keeping the
persistence boundary explicit. Money is persisted as exact decimal strings so no digit is
ever lost across a round trip~\cite{effectivejava}. The build, dependency management and
test lifecycle run on Apache Maven~\cite{maven} with JUnit~5 and MockMvc~\cite{junit5,
springtest}.

\section{Problem Statement and Research Gap}
The specific problem this project solves can be stated as follows:
\emph{Bangladeshi supershop branches lack an affordable, integrated retail management
platform that simultaneously provides (a) VAT-inclusive POS billing compliant with NBR
slab rates, (b) batch-aware inventory with expiry and shrinkage control, (c) supplier
purchasing with a guarded order lifecycle, (d) promotions and loyalty applied at the till,
(e) partial returns and voids with complete reversal, (f) a reconciled day-close Z-report,
and (g) an append-only audit trail of every sensitive action.} Existing solutions each
satisfy at most one or two of these requirements. MartFlow is the integrated response to
that research gap.

%==============================================================================
\chapter{System Analysis and Design}
%==============================================================================

\section{Technology and Tools}

\subsection{Front-End Stack}
The presentation tier is intentionally framework-free; its technologies and their roles
are summarised in Table~\ref{tab:frontend}.

\begin{table}[H]
\centering
\caption{Front-end technologies and their roles.}
\label{tab:frontend}
\small
\begin{tabularx}{\textwidth}{l l X}
\toprule
\textbf{Technology} & \textbf{Version} & \textbf{Role} \\
\midrule
HTML5 / CSS3 & --- & SPA shell (\texttt{index.html}), custom emerald theme,
thermal-receipt styling (671-line stylesheet) \\
JavaScript (ES2020+) & --- & All view logic, routing, API calls and receipt printing;
no framework by design \\
Bootstrap & 5.3.3 & Responsive layout and component vocabulary (CDN) \\
Bootstrap Icons & 1.11.3 & Icon set (CDN) \\
Fetch API / localStorage & browser-native & Authenticated API calls with the Bearer
token; token persistence across reloads \\
\bottomrule
\end{tabularx}
\end{table}

\subsection{Back-End Stack}
Table~\ref{tab:backend} lists the application-tier technologies; the Spring footprint is
deliberately the smallest workable set.

\begin{table}[H]
\centering
\caption{Back-end technologies and their roles.}
\label{tab:backend}
\small
\begin{tabularx}{\textwidth}{l l X}
\toprule
\textbf{Technology} & \textbf{Version} & \textbf{Role} \\
\midrule
Java & 17 & Language for the entire application and domain \\
Spring Boot & 3.4.0 & Embedded Tomcat, REST plumbing, JSON serialization; the only
Spring pieces used (starter-web) \\
MongoDB Java Driver & sync (BOM-managed) & Direct Atlas access through hand-written
repositories and document mappers \\
JUnit 5 + MockMvc & via starter-test & The 176-test hermetic suite \\
\bottomrule
\end{tabularx}
\end{table}

\subsection{Database}
When the \texttt{MONGODB\_URI} environment variable points at a reachable MongoDB Atlas
cluster, all data persists to Atlas (default database \texttt{martflow}); when it is absent
or unreachable, the same generic repository interface transparently serves an in-memory
implementation so the demonstration runs with zero configuration --- never half Mongo,
half memory, because all eleven repositories switch in one shot. The eleven collections
are listed in Table~\ref{tab:collections}.

\begin{table}[H]
\centering
\caption{MongoDB collections (repositories).}
\label{tab:collections}
\small
\begin{tabularx}{\textwidth}{l X}
\toprule
\textbf{Collection} & \textbf{Purpose} \\
\midrule
products & Catalog items: unit / weighed / combo, cost and MRP, VAT category, stock,
reorder level, batches with expiry \\
users & Staff accounts: PBKDF2 hash, role, active flag \\
sales & Completed bills with self-contained line snapshots, tender split, VAT, status \\
returns & Per-line returns with reasons, channels and pro-rata refund amounts \\
customers & Loyalty members: unique phone, card number, points balance \\
promotions & Category sales, member pricing and coupon codes with date windows \\
suppliers & Distributors with contact details and payment terms (Net 7/15/30) \\
purchaseOrders & POs with status, lines, goods receipts and payments \\
orderTemplates & Standing weekly restock templates (Prototype source) \\
auditLogs & Append-only activity trail (capped at 500 entries) \\
dayCloses & Z-reports: window, tender splits, counted vs expected cash, variance \\
\bottomrule
\end{tabularx}
\end{table}

\subsection{Tools Used}
Table~\ref{tab:tools} lists the supporting tools used across development, testing and
documentation.

\begin{table}[H]
\centering
\caption{Development and deployment tools.}
\label{tab:tools}
\small
\begin{tabularx}{\textwidth}{l X}
\toprule
\textbf{Tool} & \textbf{Purpose} \\
\midrule
Apache Maven (wrapper) & Build, dependency management, test lifecycle; runs with
\texttt{mvnw.cmd} and no local install \\
Visual Studio Code & Editing across Java, JavaScript, CSS and LaTeX \\
Git \& GitHub & Version control and collaboration \\
Postman & Manual API exercising during development \\
MongoDB Atlas & Free-tier cloud database and data browser \\
JUnit 5 / MockMvc & Automated hermetic testing \\
Mermaid & Diagramming for this report (architecture, DFD, ER, flows) \\
\bottomrule
\end{tabularx}
\end{table}

\section{Development Model}
The project followed an incremental model: the money path (catalog, billing, persistence)
was built first and kept runnable at all times, and management features were added as
vertical slices, each closed by tests before the next slice began. This fit a three-member
team with one shared codebase: increments were small enough to review, and the
framework-free domain let back-end and front-end work proceed in parallel against a stable
API. Table~\ref{tab:sprints} summarises the phases and Figure~\ref{fig:gantt} the
timeline.

\begin{table}[H]
\centering
\caption{Sprint plan and deliverables.}
\label{tab:sprints}
\small
\begin{tabularx}{\textwidth}{l l X}
\toprule
\textbf{Phase} & \textbf{Duration} & \textbf{Deliverable} \\
\midrule
Planning \& Design      & Weeks 1--2   & Requirements, 18-pattern map, ER model,
wireframes \\
Domain \& Back-End      & Weeks 3--5   & Pattern-bearing domain, persistence +
fallback, auth, 15 controllers / 62 endpoints \\
Front-End               & Weeks 6--9   & SPA shell, router + role gate, POS screen,
management screens \\
Feature Enhancement     & Weeks 10--11 & Day-close, audit trail, promotions manager,
sales explorer, Developer Mode \\
Testing \& QA           & Weeks 12--13 & 176 hermetic tests, restart drills,
drawer-math checks, hardening \\
Documentation           & Weeks 14--15 & README, Bengali companion doc, drift guards,
this report, presentation \\
\bottomrule
\end{tabularx}
\end{table}

\begin{figure}[H]
\centering
\includegraphics[width=\textwidth]{gantt-sprints.png}
\caption{Fifteen-week sprint plan Gantt chart.}
\label{fig:gantt}
\end{figure}

\section{System Architecture}
MartFlow is a three-tier web application (Figure~\ref{fig:arch}). The browser runs the
vanilla-JavaScript SPA; every request carries the Bearer token through the
\texttt{AuthFilter}, which resolves the caller into a ThreadLocal role context. The 15
REST controllers are deliberately thin --- each handler is essentially one facade call ---
so HTTP concerns never leak into business logic. The facades (\texttt{MartFlowFacade},
\texttt{BillingFacade}, \texttt{PurchasingService} and siblings) orchestrate the plain-Java
domain packages where the 18 design patterns live. At the data boundary a protection proxy
re-checks role policy on product writes before the generic repository interface forwards
to either MongoDB Atlas or the in-memory fallback. Spring wires this entire object graph
in a single configuration class of 11 bean methods; the domain itself contains no
framework annotation at all.

\begin{figure}[H]
\centering
\includegraphics[width=0.98\textwidth]{system-architecture.png}
\caption{Three-tier system architecture of MartFlow.}
\label{fig:arch}
\end{figure}

\section{GoF Design Patterns Applied}
\label{sec:patterns}
Because this project is, by course, a design-pattern project, the pattern catalogue is
treated as part of the system design rather than an implementation accident. The governing
rule was \emph{a pattern earns its place only by carrying a real feature}. Exactly 18
patterns are applied --- 4~creational, 5~structural and 9~behavioral --- a split asserted
by an automated test, and each pattern is bound to the business reason it exists, as
Table~\ref{tab:patterns} documents.

\begin{longtable}{p{2.9cm}p{4.3cm}p{7.2cm}}
\caption{The 18 GoF design patterns applied in MartFlow.}
\label{tab:patterns}\\
\toprule
\textbf{Pattern} & \textbf{Where} & \textbf{Business reason} \\
\midrule
\endfirsthead
\caption[]{(continued) The 18 GoF design patterns applied in MartFlow.}\\
\toprule
\textbf{Pattern} & \textbf{Where} & \textbf{Business reason} \\
\midrule
\endhead
\bottomrule
\endfoot
\bottomrule
\endlastfoot
Singleton & \texttt{InventoryCatalog}, \texttt{DatabaseConnection} & One stock truth
per process; two catalogs would oversell the same shelf. \\
Factory Method & \texttt{ProductFactory} with \texttt{Unit}- and
\texttt{WeighedProductFactory} & The add-item form creates per-piece and per-kg items
through one path; no type \texttt{if/else}. \\
Builder & \texttt{PurchaseOrderBuilder} & A PO is assembled line by line with
validation; no half-configured order can escape into the workflow. \\
Prototype & \texttt{StandingOrderTemplate} & The weekly restock list is cloned into a
draft --- two clicks instead of fifteen. \\
Adapter & \texttt{payment/} cash, card, bKash, Nagad, points adapters over vendor-shaped
SDKs & Five very different tender channels behind one port; the till never branches on
payment type. \\
Decorator & \texttt{LineDiscount}, \texttt{CarryBagFee}, \texttt{DeliveryFee},
\texttt{RoundOffAdjustment} & Charges and promotions stack on any bill line without
touching the line classes. \\
Facade & \texttt{MartFlowFacade}, \texttt{BillingFacade} & Every UI button is one call;
controllers stay one line and testable. \\
Proxy (protection) & \texttt{RoleGuardProxy} over the repositories & A cashier cannot
reprice the catalog even through a buggy endpoint --- enforced at the data boundary. \\
Composite & \texttt{ComboProduct} & An Eid hamper is priced and stocked from its
components; selling one consumes each component truthfully (stock = the scarcest). \\
Observer & \texttt{AlertService}, \texttt{ReorderSuggestionObserver},
\texttt{StockSubject} & Low stock, near-expiry batches and wastage raise alerts by
themselves; low stock even drafts the reorder suggestion. \\
Strategy & \texttt{RegularPrice}, \texttt{CategorySale}, \texttt{MemberPrice} +
\texttt{PromotionEngine} & Managers toggle promotions per date window; the till's pricing
code never changes. \\
Command & \texttt{billing/commands/} + void pipeline & Tender is a transaction: a
declined card rolls back stock, points and the sale itself in reverse order. \\
Template Method & \texttt{AbstractReportGenerator} + 8 concrete reports & A new report
is three hooks; ``VOIDED never counts as revenue'' is decided exactly once. \\
Chain of Responsibility & \texttt{billing/validation/} --- 6 ordered rules & The cashier
sees the exact first failing rule (empty bill, weighment, stock, coupon, loyalty,
tender). \\
State & \texttt{suppliers/postate/} --- 6 purchase-order states & A draft PO physically
cannot be received; a closed one cannot be cancelled. \\
Visitor & \texttt{VatVisitor}, \texttt{ProfitVisitor},
\texttt{ReceiptFormatterVisitor} & The same bill lines feed VAT filing, margin analysis
and receipt layout --- including sales reloaded from the database. \\
Memento & \texttt{BillMemento} (per-session undo stack, depth 10) & Mis-scans happen
every minute at a real till; the cashier's Undo is a money feature. \\
Iterator & \texttt{catalog/iter/} --- in-stock, low-stock, expiring, price-range & The
low-stock view and the low-stock report walk the same iterator --- they can never
disagree. \\
\end{longtable}

Two deliberate asymmetries are part of the design answer. A sale's status is a plain enum
while purchase orders get the full State pattern, because a sale has no long-lived
per-status behaviour (void and return are one-shot operations), whereas a PO sits for days
in states with genuinely different allowed operations. And several pattern candidates
were rejected on the merits: the audit trail does not use Observer (stock events carry no
actor or intent), the Z-report does not use the report Template Method (its
``VOIDED never counts'' invariant is exactly what a Z-report must break for cash
purposes), and no nineteenth pattern was added for symmetry. The rejection rationale is
written into each service's documentation --- the ``why not'' is part of the answer.

\section{Use-Case Analysis}
Four actors interact with the system. The cashier runs the till; the manager (who
inherits every cashier use case) runs the floor; the owner (who inherits every manager
use case) additionally controls the catalog's destruction and the staff accounts; and a
developer account exists solely for examination --- it sees the till-side screens and the
Developer Mode, but holds no business power. Table~\ref{tab:usecases} summarises the
principal use cases and Figure~\ref{fig:usecase} shows the diagram.

\begin{table}[H]
\centering
\caption{Use-case summary by actor.}
\label{tab:usecases}
\small
\begin{tabularx}{\textwidth}{l X}
\toprule
\textbf{Actor} & \textbf{Representative use cases} \\
\midrule
Cashier & Log in; scan and bill items (unit, weighed, combo); attach loyalty member and
coupon; split tender across five channels; undo mis-scans; print receipt; process
returns; register loyalty member; view operational reports; read alerts \\
\addlinespace
Manager & Inherits all cashier use cases; manage catalog items (add, edit, restock,
adjust); manage promotions and test coupons; manage suppliers, purchase orders, goods
receipts, payments and restock templates; void sales with a reason; run financial
reports; close the day (Z-report); inspect the activity log \\
\addlinespace
Owner (admin) & Inherits all manager use cases; delete catalog items; manage staff
accounts (create, disable, reset password) \\
\addlinespace
Developer & Till-side screens as the cashier; open Developer Mode (pattern catalogue,
API explorer, diagnostics) --- and nothing else \\
\bottomrule
\end{tabularx}
\end{table}

\begin{figure}[H]
\centering
\includegraphics[width=\textwidth]{use-case.png}
\caption{System Use Case Diagram.}
\label{fig:usecase}
\end{figure}

\section{Context-Level Diagram (DFD-0)}
The context diagram (Figure~\ref{fig:dfd0}) shows MartFlow as a single process exchanging
flows with its external entities: the three staff roles, the payment vendors whose
SDK shapes sit behind the adapters, and the MongoDB Atlas data store.

\begin{figure}[H]
\centering
\includegraphics[width=0.98\textwidth]{dfd-0.png}
\caption{Context-level diagram (DFD-0).}
\label{fig:dfd0}
\end{figure}

\section{Data-Flow Diagram (DFD-1)}
The first-level decomposition (Figure~\ref{fig:dfd1}) breaks the system into eight
cooperating processes --- authentication and access, catalog and inventory, POS billing
and tender, returns and voids, purchasing, promotions and loyalty, reports and day close,
and alerts with audit --- all of which read from and write to the shared data stores.

\begin{figure}[H]
\centering
\includegraphics[width=\textwidth]{dfd-1.png}
\caption{First-level data-flow diagram (DFD-1).}
\label{fig:dfd1}
\end{figure}

\section{Database Schema}
The entity relationships among the eleven collections are shown in
Figure~\ref{fig:erdiagram}. Money is stored as exact decimal strings and sale lines are
self-contained snapshots (name, SKU, price, VAT rate and cost frozen at sale time), which
is what lets a receipt reloaded from the database reconstruct into the identical item
chain and print byte-identically. Selected field-level detail is given in
Table~\ref{tab:schema}.

\begin{figure}[H]
\centering
\includegraphics[width=\textwidth]{er-diagram.png}
\caption{Entity-relationship overview of the MartFlow database.}
\label{fig:erdiagram}
\end{figure}

\begin{table}[H]
\centering
\caption{Selected schema field detail.}
\label{tab:schema}
\small
\begin{tabularx}{\textwidth}{l l X}
\toprule
\textbf{Model} & \textbf{Key field} & \textbf{Notes} \\
\midrule
Product & \texttt{type} & Enum \texttt{UNIT} / \texttt{WEIGHED} / \texttt{COMBO};
batches embedded with batch no and expiry; VAT rate resolved through the category \\
Sale & \texttt{status} & Enum \texttt{COMPLETED} / \texttt{PARTIALLY\_RETURNED} /
\texttt{RETURNED} / \texttt{VOIDED}; lines are self-contained snapshots; money stored as
exact decimal strings \\
SaleReturn & \texttt{lines} & Per-line quantity and reason; pro-rata refund amounts;
channel defaults to the original tender \\
User & \texttt{passwordHash} & \texttt{base64(salt):base64(key)} PBKDF2 format; role
enum; \texttt{active} flag for disable instead of delete \\
PurchaseOrder & \texttt{status} & Six states guarded by the State pattern
(\texttt{DRAFT} \textrightarrow{} \texttt{ORDERED} \textrightarrow{}
\texttt{PARTIALLY\_RECEIVED} \textrightarrow{} \texttt{RECEIVED} \textrightarrow{}
\texttt{CLOSED}, plus \texttt{CANCELLED}); receipts and payments embedded \\
Promotion & \texttt{type} & Enum \texttt{CATEGORY\_SALE} / \texttt{MEMBER\_PRICE} /
\texttt{COUPON\_FLAT} / \texttt{COUPON\_PERCENT}; active date window \\
Customer & \texttt{phone} & Unique key for till lookup; integer points balance \\
AuditLog & \texttt{action} & Append-only; capped at 500 entries; actor always recorded \\
DayClose & \texttt{variance} & \texttt{countedCash} minus \texttt{expectedCash}; zero
renders as balanced \\
\bottomrule
\end{tabularx}
\end{table}

\section{Sequence Diagram}
The sequence diagram (Figure~\ref{fig:seqdiagram}) illustrates the execution flow for
the system's money path --- tendering a bill. The validation chain runs first and stops
the request on the exact first failing rule; the command pipeline then reserves stock,
charges the tenders, awards points and persists the sale; and a declined card triggers a
last-in-first-out rollback so that, from the shop's point of view, nothing happened.

\begin{figure}[H]
\centering
\includegraphics[width=0.98\textwidth]{sequence-pos-sale.png}
\caption{Sequence Diagram for POS tender with rollback.}
\label{fig:seqdiagram}
\end{figure}

\section{Activity Diagram}
The activity diagram (Figure~\ref{fig:actdiagram}) details the cashier's primary journey
--- building a bill and taking payment --- including the weighment and validation loops
the till enforces.

\begin{figure}[H]
\centering
\includegraphics[width=0.82\textwidth]{activity-billing.png}
\caption{Activity Diagram for the POS billing journey.}
\label{fig:actdiagram}
\end{figure}

\section{State Transition: Purchase Order Lifecycle}
Purchase orders are the system's clearest stateful entity, and the State pattern guards
their six states (Figure~\ref{fig:postate}): a draft physically cannot be received, and a
closed order cannot be cancelled --- the states refuse illegal operations themselves
rather than relying on the caller's discipline.

\begin{figure}[H]
\centering
\includegraphics[width=0.72\textwidth]{state-po-lifecycle.png}
\caption{Purchase-order state machine (State pattern).}
\label{fig:postate}
\end{figure}

\section{Algorithms and Flowcharts}

\subsection{VAT Back-Calculation per NBR Slabs}
Shelf prices in Bangladesh are VAT-inclusive, so output VAT must be derived, not added.
For a line with net amount $n$ and category rate $r \in \{0, 7.5, 15\}$ percent, the
line's output VAT is
\begin{equation}
\mathrm{VAT} = n \times \frac{r}{100 + r},
\end{equation}
computed in \texttt{BigDecimal} at scale~2 with \texttt{HALF\_UP} rounding, per line
first and summed afterwards. Worked examples pinned by tests: 115.00 at 15\% gives
15.00; 107.50 at 7.5\% gives 7.50; 99.00 at 0\% gives 0.00. Figure~\ref{fig:vatflow}
shows the flow.

\begin{figure}[H]
\centering
\includegraphics[width=0.72\textwidth]{flow-vat-calculation.png}
\caption{VAT back-calculation flowchart (NBR slabs).}
\label{fig:vatflow}
\end{figure}

\subsection{Tender Command Pipeline with Rollback}
Tender is modelled as a command transaction: four commands run in order
(reserve stock, charge tenders, award points, create the sale), and any exception
triggers a last-in-first-out, best-effort undo of the commands already executed --- each
undo wrapped in its own \texttt{try/catch} so one failure cannot mask another. A
declined card therefore leaves no reserved stock, no awarded points and no sale; the API
answers 409 and the cashier sees exactly what went wrong. The same three-command shape
(reserve, refund, mark voided) reverses a voided sale. Figure~\ref{fig:tenderflow}
shows the flow.

\begin{figure}[H]
\centering
\includegraphics[width=0.78\textwidth]{flow-tender-rollback.png}
\caption{Tender command pipeline with LIFO rollback.}
\label{fig:tenderflow}
\end{figure}

\subsection{Authentication Flow}
Every staff member authenticates the same way: the SPA posts credentials to
\texttt{POST /api/auth/login}; the server verifies the PBKDF2 hash (running a dummy
verification on unknown usernames so timing cannot reveal which accounts exist) and
issues a 256-bit random hex token with a sliding 12-hour lifetime. Every subsequent
request carries the token as a Bearer header; the \texttt{AuthFilter} resolves it to a
caller and stores it in a ThreadLocal role context that is cleared in a
\texttt{finally} block, after which controller-level role gates and the data-boundary
proxy enforce policy. Figure~\ref{fig:authflow} shows the flow.

\begin{figure}[H]
\centering
\includegraphics[width=0.85\textwidth]{auth-flow.png}
\caption{Authentication and authorization flow.}
\label{fig:authflow}
\end{figure}

%==============================================================================
\chapter{Implementation}
%==============================================================================

\section{Front-End / Interface Design}
The user interface is a framework-free single-page application of roughly 2{,}800 lines
of JavaScript across eight modules, served by the back end from
\texttt{src/main/resources/static/}. One HTML shell hosts a view host that the router
fills; every screen is a render function registered against a route.

\paragraph{Design system.} Bootstrap~5.3.3 supplies the layout and component vocabulary,
restyled by a 671-line custom stylesheet around an emerald and gold palette; Bootstrap
Icons~1.11.3 provides the iconography. A dark/light theme toggle is persisted per
browser, and the thermal receipt is rendered in a fixed 38-column style that matches a
320\,px slip.

\paragraph{Routing and access control.} A 67-line hash router registers the 15 routes
(each with a title, icon and role list). On every hash change it clears the previous
view's cleanup callback, checks authentication, and enforces the role list --- a cashier
navigating to a manager screen receives ``Your role cannot open that page'' and is
redirected to the POS. The navigation bar itself is generated from the same route table,
so menus and gates can never disagree.

\paragraph{API layer.} A single \texttt{api.js} module wraps \texttt{fetch}: it attaches
the Bearer token from \texttt{localStorage}, JSON-encodes bodies, unwraps the error
contract (400 validation, 401 session-expired, 403 access denied, 409 conflict), and
raises uniform toasts so individual views never duplicate error handling.

\paragraph{Views.} Screen logic is split across three view modules: till-side screens
(POS, receipt), operations (inventory, purchasing, loyalty, returns, alerts) and
management (dashboard, sales explorer, promotions, reports, day close, activity log,
staff, developer mode). The POS screen provides the scan box with live suggestions,
category chips, a 24-tile quick-pick grid, per-line edit with a 0.25-step weight control,
and an undo button that reflects the server-side memento depth.

\paragraph{Receipt printing.} A dedicated receipt module renders the slip, supports
thermal printing through a hidden iframe, copying the text, and saving it as a plain-text
file; pressing Enter after a sale arms the next customer automatically.

\section{Back-End}

\paragraph{Entry point and wiring.} \texttt{MartFlowApplication} boots Spring with the
Mongo auto-configuration excluded (the raw driver owns its client), and
\texttt{AppConfig} is the single configuration class: eleven bean methods assemble the
entire object graph as plain objects --- catalogs over a role-guarded repository,
observers subscribed to stock events, facades receiving payment adapters and the
promotion engine --- and run the seed data on first boot (37 products, 9 categories,
4 promotions, 5 loyalty members, 3 suppliers and an in-flight purchase order).

\paragraph{Request pipeline.} The \texttt{AuthFilter} validates the Bearer token against
the token store, places the caller in the ThreadLocal \texttt{RoleContext}, and always
clears it in a \texttt{finally} block (Listing~\ref{lst:rolecontext}); controllers then
enforce policy with role gates before delegating one line to a facade. A global
exception handler maps domain errors onto the HTTP contract.

\begin{lstlisting}[language=Java,caption={Request context that never leaks between
requests (simplified).},label={lst:rolecontext}]
// AuthFilter, simplified
try {
    Caller caller = tokenStore.resolve(bearerToken);
    RoleContext.set(caller);
    chain.doFilter(request, response);
} finally {
    RoleContext.clear();  // no cross-request role leakage
}
\end{lstlisting}

\paragraph{Controllers.} The 15 REST controllers expose 62 handler endpoints under
\texttt{/api}; each is a thin translation between JSON and the facades, so business rules
exist exactly once, in the domain. Selected endpoints are tabulated in
Appendix~\ref{app:endpoints}.

\paragraph{The money rule.} Every amount is a \texttt{BigDecimal} of scale~2 with
\texttt{HALF\_UP} rounding produced through one \texttt{MoneyUtil} choke point;
quantities carry scale~3 for weighed goods. Cash tenders round the payable to whole
taka through a decorator that shows the adjustment on the receipt; card and mobile-wallet
tenders do not round. Listing~\ref{lst:vat} shows the shared VAT calculation.

\begin{lstlisting}[language=Java,caption={Inclusive-VAT back-calculation, the single
source of NBR output VAT (simplified).},label={lst:vat}]
// VatCalculator, simplified
public BigDecimal vatOf(String categoryId, BigDecimal net) {
    BigDecimal rate = categories.vatRate(categoryId); // 0, 7.5 or 15
    if (rate.signum() == 0) return BigDecimal.ZERO;
    return net.multiply(rate)
              .divide(rate.add(HUNDRED), 2, RoundingMode.HALF_UP);
}
\end{lstlisting}

\paragraph{Security design.} Authorization is enforced in three deliberate layers: the
menu hides what a role cannot do (cosmetic), the controller role gate rejects the call
(real), and the protection proxy re-checks policy on product writes at the data boundary
(last line). The client never declares its role; it always derives from the token.
Password hashing follows RFC~8018 (PBKDF2, 120{,}000 iterations, per-user salt,
constant-time comparison), and unknown usernames run a dummy verification so response
timing cannot enumerate accounts.

\paragraph{Persistence.} A static selector probes the Atlas connection once at startup
and wires all eleven repositories to either the Mongo implementations (with hand-written
document mappers) or the in-memory fallbacks --- never a mixture. Money is stored as
exact decimal strings, and sale lines are self-contained snapshots: a receipt reloaded
from the database reconstructs the identical decorator chain and prints byte-identically,
which is what keeps reprints and VAT filings from drifting from what was charged.

\paragraph{Developer Mode and drift guards.} The developer-only endpoints serve an
in-product catalogue of all 18 patterns (each card cites the real class, a real source
snippet and the test that proves it), a live API explorer that fires any route with the
caller's own token --- a 403 demonstrates the gate --- and system diagnostics. Two
automated guards keep this documentation honest: one test walks the source tree and
fails the build if a pattern card cites a class that does not exist, and a second
asserts the API explorer's catalogue equals Spring's registered routes in both
directions.

\section{Modules and Features}
The following 16 functional modules were implemented:

\begin{enumerate}
  \item \textbf{Authentication and Access Control} --- PBKDF2 login, bearer sessions,
        four roles enforced from menu to data boundary, failed-login auditing.
  \item \textbf{POS Billing} --- scan, search and quick-pick entry; weighed goods;
        combos; per-line editing with undo (depth 10); customer, coupon and charge
        attachment; split tender across five channels; thermal receipts.
  \item \textbf{Inventory Management} --- cost/MRP per item, batches with expiry,
        reorder levels, low-stock and expiring views, damage/loss/theft/count
        adjustments with shrinkage alerts.
  \item \textbf{Purchasing and PO Lifecycle} --- suppliers with payment terms, guarded
        six-state purchase orders, goods receipts that land stock with batch and updated
        cost, payment records, payables.
  \item \textbf{Standing Restock Templates} --- weekly order lists saved as templates
        and cloned into fresh drafts (two clicks instead of fifteen).
  \item \textbf{Promotions Engine} --- date-windowed category sales, member pricing and
        coupon codes with manager CRUD, a power toggle that reprices the till instantly,
        and a live coupon tester.
  \item \textbf{Loyalty Program} --- members by unique phone, points earned per 100 BDT
        net, points redeemed as tender (1 point = 1 BDT), manager adjustments.
  \item \textbf{Returns and Exchanges} --- partial per-line returns with pro-rata refunds
        through the original channel, repeat-return protection, exchange as return plus
        new bill.
  \item \textbf{Sales Explorer and Void} --- every receipt with date/status/cashier
        filters, one-click reprint, and manager void with a mandatory reason and full
        reversal.
  \item \textbf{Day-Close Z-Report} --- per-tender split, expected versus counted cash
        with variance, printable Z-slip, history reprint; a sale voided today charges
        today's drawer even if sold yesterday.
  \item \textbf{Alerts} --- automatic low-stock, expiry, shrinkage, restock and
        price-change alerts with a suggestion service, unread badge and mark-read.
  \item \textbf{Activity Log} --- append-only audit of logins (including failures),
        voids, returns, shrinkage, price edits, promotion changes, PO moves and staff
        changes, always with the actor.
  \item \textbf{Reports and Dashboard} --- eight report types (daily sales, best
        sellers, low-stock worksheet, expiry watch, returns, staff performance, profit,
        NBR VAT summary) each exportable to CSV, plus a six-KPI operations dashboard.
  \item \textbf{Staff Management} --- owner-only account creation with role, disabling
        and password reset.
  \item \textbf{Developer Mode} --- pattern deep-dives, live API explorer and system
        diagnostics, visible only to the developer account.
  \item \textbf{Developer Diagnostics} --- persistence mode, active token, repository
        counts, uptime and the Dhaka clock for examination and support.
\end{enumerate}

The dashboard's tender distribution over a seeded demonstration session is illustrated
in Figure~\ref{fig:pietender}.

\begin{figure}[H]
\centering
\includegraphics[width=0.62\textwidth]{pie-tender-mix.png}
\caption{Tender mix on the demonstration dashboard (illustrative data).}
\label{fig:pietender}
\end{figure}

%==============================================================================
\chapter{User Manual}
%==============================================================================

\section{System Requirements}

\subsection{Hardware Requirements}
Table~\ref{tab:hw} lists the minimum and recommended hardware for running the suite.

\begin{table}[H]
\centering
\caption{Hardware requirements.}
\label{tab:hw}
\small
\begin{tabularx}{\textwidth}{l X X}
\toprule
\textbf{Component} & \textbf{Minimum} & \textbf{Recommended} \\
\midrule
Processor & Intel Core i3 / AMD Ryzen 3 & Core i5 / Ryzen 5 or better \\
RAM & 4\,GB & 8\,GB \\
Storage & 1\,GB free (app + logs) & 5\,GB free (Atlas is cloud-hosted) \\
Display & 1366$\times$768 & 1920$\times$1080 \\
Peripherals & Keyboard; barcode scanner (optional) & Scanner + thermal printer \\
Network & Broadband for Atlas and the CDN & Stable broadband \\
\bottomrule
\end{tabularx}
\end{table}

\subsection{Software Requirements}
\paragraph{For users.} Any modern desktop or tablet browser (Chrome, Edge, Firefox or
Safari, current versions) with JavaScript enabled and internet access for the Bootstrap
CDN; no installation is needed on the till machine beyond the browser.
\paragraph{For local development.} Windows 10/11 (or any OS), JDK~17 or newer, the
included Maven wrapper (no separate Maven install), optionally a free MongoDB Atlas
account with a connection string, Visual Studio Code and Git.

\section{User Interfaces}

\subsection{Panel A --- Till-Side Screens}
The cashier logs in and lands on the POS (Figure~\ref{fig:pos}): a scan box with live
suggestions, category chips and a quick-pick grid build the bill; per-line edit, undo and
clear sit beside the totals ladder (Gross, Promotions, Coupon, charges, PAYABLE including
VAT). Tendering opens the split-payment modal across cash, card, bKash, Nagad and points,
and completes into the receipt (Figure~\ref{fig:receipt}) with its VAT breakdown, round-off
line, transaction references and change; Enter returns the till to the next customer.

\uifig{The POS billing screen.}{fig:pos}{figs/fig5_1_pos.png}
\uifig{The printed thermal receipt.}{fig:receipt}{figs/fig5_2_receipt.png}

\subsection{Panel B --- Operations Screens}
Inventory (Figure~\ref{fig:inventory}) lists items with type badges, cost, MRP,
stock highlighted when low, reorder levels and batch expiries; managers restock, adjust
shrinkage and edit prices from here. Purchasing (Figure~\ref{fig:purchases}) hosts the
purchase-order board with its state badges, the goods-receipt dialog that lands stock
with batch and cost, supplier records with payment terms, and the standing templates that
clone into weekly drafts. Promotions (Figure~\ref{fig:promotions}) manages the four
promotion types with a live coupon tester, and returns reconcile receipts per line with
pro-rata refunds.

\uifig{Inventory with batches, reorder levels and shrinkage controls.}{fig:inventory}{figs/fig5_3_inventory.png}
\uifig{Purchasing: PO board, goods receipts, suppliers and templates.}{fig:purchases}{figs/fig5_4_purchases.png}
\uifig{Promotions manager with the live coupon tester.}{fig:promotions}{figs/fig5_7_promotions.png}

\subsection{Panel C --- Management Screens}
The dashboard (Figure~\ref{fig:dashboard}) shows the six operational KPIs --- net sales,
bill count, average basket, output VAT, cash in drawer, low-stock and expiring counts ---
with quick report links and the alert feed. The sales explorer filters every receipt and
reprints or voids with a reason; day close (Figure~\ref{fig:dayclose}) previews the
window's tender split, takes the counted drawer and prints the Z-slip with the variance;
the activity log (Figure~\ref{fig:activity}) is the searchable audit trail; and Developer
Mode (Figure~\ref{fig:developer}) exposes the 18 pattern cards, the API explorer and
diagnostics to the examiner's account only.

\uifig{The operations dashboard.}{fig:dashboard}{figs/fig5_5_dashboard.png}
\uifig{Day-close Z-report with drawer variance.}{fig:dayclose}{figs/fig5_8_dayclose.png}
\uifig{The append-only activity log.}{fig:activity}{figs/fig5_9_activity.png}
\uifig{Developer Mode: pattern deep-dives with real source and tests.}{fig:developer}{figs/fig5_6_developer.png}

\subsection{Login Credentials}
The following demonstration accounts are seeded on first boot (change them before any
real use):
\begin{description}[leftmargin=1.5em]
  \item[\textbf{Owner (admin):}] full access including staff management and item
        deletion
        \begin{itemize}
          \item \textbf{Username:} \texttt{admin}
          \item \textbf{Password:} \texttt{admin123}
        \end{itemize}
  \item[\textbf{Manager:}] catalog, purchasing, financial reports, voids, day close
        \begin{itemize}
          \item \textbf{Username:} \texttt{manager}
          \item \textbf{Password:} \texttt{manager123}
        \end{itemize}
  \item[\textbf{Cashier:}] billing, returns, operational views
        \begin{itemize}
          \item \textbf{Username:} \texttt{cashier}
          \item \textbf{Password:} \texttt{cashier123}
        \end{itemize}
  \item[\textbf{Developer:}] till-side screens plus Developer Mode only
        \begin{itemize}
          \item \textbf{Username:} \texttt{developer}
          \item \textbf{Password:} \texttt{developer123}
        \end{itemize}
\end{description}

\section{Getting Started}
\begin{enumerate}
  \item Install JDK~17 or newer and clone the repository.
  \item Run the application with the included wrapper (no Maven install required):
        \texttt{.\textbackslash mvnw.cmd spring-boot:run} --- the suite starts in
        in-memory mode with seeded data and prints the demo logins.
  \item Open \url{http://localhost:8080} and sign in with a demo account.
  \item Optional, for durable storage: create a free MongoDB Atlas cluster, copy the
        connection string, and set \texttt{MONGODB\_URI} (and optionally
        \texttt{MONGODB\_DB}) as an environment variable or in a \texttt{.env} file ---
        see \texttt{.env.example}. When Atlas is reachable, sales, stock, points and
        purchase orders survive restarts.
  \item Run the test suite any time with \texttt{.\textbackslash mvnw.cmd clean test}
        (176 tests; the suite pins an invalid database URI so it can never touch a real
        database).
\end{enumerate}

%==============================================================================
\chapter{Conclusion}
%==============================================================================

\section{Conclusion}
MartFlow matured into a complete supershop management suite: a vanilla-JavaScript SPA of
15 role-gated routes over a Java~17 / Spring Boot~3.4.0 application layer of 15 REST
controllers and 62 endpoints, a framework-free domain in which exactly 18 GoF design
patterns (4~creational, 5~structural, 9~behavioral) each carry a real business feature,
and a MongoDB Atlas data tier of 11 collections with a zero-configuration in-memory
fallback. The money path is honest by construction --- VAT-inclusive NBR slab
back-calculation, exact-decimal arithmetic, command transactions with rollback, and a
day-close that reconciles by hand --- and the accountability layer (audit trail plus
Z-report) answers ``who did what'' at any point. Security is enforced in three layers
across four roles down to the data boundary, and 176 automated tests in 31 hermetic
classes, together with the in-product Developer Mode and its drift guards, keep both the
code and its documentation provable. The system is deployable today as a single
executable JAR for a single branch, with persistence on Atlas or nothing at all.

\section{Limitations}
\begin{enumerate}
  \item \textbf{Simulated vendor SDKs.} The card, bKash and Nagad integrations run
        against vendor-shaped fakes; real merchant credentials would change only the ten
        to fifteen lines inside each adapter, but are out of reach of a course project.
  \item \textbf{CDN dependency.} Bootstrap loads from a CDN, so a fully offline till
        needs pre-cached assets.
  \item \textbf{In-memory sessions.} Login tokens and the alert feed are deliberately
        ephemeral; a restart asks staff to sign in again (business data is never lost).
  \item \textbf{No concurrency stress test.} Race defenses exist at three levels
        (per-token bill sessions, synchronized stock movement, ThreadLocal contexts) and
        are covered by single-request tests, but a dedicated multi-thread stress suite is
        future work.
  \item \textbf{Single branch.} One deployment serves one shop; multi-branch
        synchronisation is not yet modelled.
  \item \textbf{Online-only till.} Without connectivity to the server the till cannot
        bill; offline-first operation with background reconciliation is future work.
  \item \textbf{Illustrative analytics.} Dashboard aggregates and chart figures in this
        report reflect seeded demonstration data rather than long-run production
        history.
\end{enumerate}

\section{Future Works}
The natural next steps follow the roadmap set during development: multi-branch
synchronisation under one owner; offline-first billing with background reconciliation;
real bKash, Nagad and card-acquirer integrations behind the existing adapter port;
barcode-printer and cash-drawer hardware support; customer purchase-history analytics
over the sale snapshots; and an approval workflow for purchase orders. On the engineering
side: persisting sessions across restarts, a CountDownLatch-based concurrency stress
suite, and an automated restart drill in continuous integration.

%==============================================================================
%  BACK MATTER
%==============================================================================
\begin{thebibliography}{99}
\addcontentsline{toc}{chapter}{References}

\bibitem{java17}     Oracle. ``Java SE 17 Documentation.'' Oracle Corporation,
                     \url{https://docs.oracle.com/en/java/javase/17/}, 2026.
\bibitem{springboot} VMware Tanzu. ``Spring Boot Reference Documentation, v3.4.0.''
                     \url{https://docs.spring.io/spring-boot/3.4/}, 2026.
\bibitem{spring}    VMware Tanzu. ``Spring Framework Documentation.''
                     \url{https://docs.spring.io/spring-framework/reference/}, 2026.
\bibitem{springtest} VMware Tanzu. ``Testing the Web Layer (MockMvc).'' Spring.io
                     Guides, \url{https://spring.io/guides/gs/testing-web}, 2026.
\bibitem{mongodb}   MongoDB, Inc. ``MongoDB Manual.''
                     \url{https://www.mongodb.com/docs/manual/}, 2026.
\bibitem{atlas}     MongoDB, Inc. ``MongoDB Atlas Documentation.''
                     \url{https://www.mongodb.com/docs/atlas/}, 2026.
\bibitem{mongodriver} MongoDB, Inc. ``MongoDB Java Driver (sync) Documentation.''
                     \url{https://www.mongodb.com/docs/drivers/java/sync-current/}, 2026.
\bibitem{maven}     Apache Software Foundation. ``Apache Maven Documentation.''
                     \url{https://maven.apache.org/guides/}, 2026.
\bibitem{junit5}    JUnit Team. ``JUnit 5 User Guide.''
                     \url{https://junit.org/junit5/docs/current/user-guide/}, 2026.
\bibitem{gof}       E.~Gamma, R.~Helm, R.~Johnson and J.~Vlissides, \emph{Design
                     Patterns: Elements of Reusable Object-Oriented Software}.
                     Reading, MA: Addison-Wesley, 1994.
\bibitem{headfirst} E.~Freeman and E.~Robson, \emph{Head First Design Patterns},
                     2nd~ed. Sebastopol, CA: O'Reilly Media, 2021.
\bibitem{refactoringguru} Refactoring.Guru. ``Design Patterns.''
                     \url{https://refactoring.guru/design-patterns}, 2026.
\bibitem{effectivejava} J.~Bloch, \emph{Effective Java}, 3rd~ed. Boston, MA:
                     Addison-Wesley, 2018 (Item~60: ``Avoid float and double if exact
                     answers are required'').
\bibitem{cleanarch} R.~C.~Martin, \emph{Clean Architecture: A Craftsman's Guide to
                     Software Structure and Design}. Boston, MA: Prentice Hall, 2017.
\bibitem{sommerville} I.~Sommerville, \emph{Software Engineering}, 10th~ed. Boston,
                     MA: Pearson, 2016.
\bibitem{fielding}  R.~T.~Fielding, ``Architectural Styles and the Design of
                     Network-based Software Architectures,'' Ph.D. dissertation,
                     University of California, Irvine, 2000.
\bibitem{rfc8018}   K.~Moriarty, Ed., ``PKCS \#5: Password-Based Cryptography
                     Specification Version 2.1,'' RFC~8018, Internet Research Task
                     Force, 2017.
\bibitem{owasp}     OWASP Foundation. ``Password Storage Cheat Sheet.''
                     \url{https://cheatsheetseries.owasp.org/cheatsheets/Password_Storage_Cheat_Sheet.html},
                     2026.
\bibitem{nbr}       National Board of Revenue, Government of Bangladesh. ``Value Added
                     Tax and Supplementary Duty Act 2012 (as amended).'' Dhaka: NBR,
                     2012.
\bibitem{bootstrap} Bootstrap Team. ``Bootstrap 5.3 Documentation.''
                     \url{https://getbootstrap.com/docs/5.3/}, 2026.
\bibitem{bootstrapicons} Bootstrap Team. ``Bootstrap Icons.''
                     \url{https://icons.getbootstrap.com/}, 2026.
\bibitem{mdnjs}     Mozilla. ``JavaScript --- MDN Web Docs.''
                     \url{https://developer.mozilla.org/en-US/docs/Web/JavaScript}, 2026.
\bibitem{mdnfetch}  Mozilla. ``Fetch API --- MDN Web Docs.''
                     \url{https://developer.mozilla.org/en-US/docs/Web/API/Fetch_API}, 2026.
\bibitem{mdnlocal}  Mozilla. ``Window: localStorage --- MDN Web Docs.''
                     \url{https://developer.mozilla.org/en-US/docs/Web/API/Window/localStorage},
                     2026.
\bibitem{git}       Software Freedom Conservancy. ``Git Documentation.''
                     \url{https://git-scm.com/doc}, 2026.
\bibitem{postman}   Postman, Inc. ``Postman Documentation.''
                     \url{https://learning.postman.com/docs/}, 2026.
\bibitem{mermaid}   Mermaid. ``Mermaid Diagramming and Charting Documentation.''
                     \url{https://mermaid.js.org/intro/}, 2026.
\bibitem{dfd}       Lucidchart. ``What is a Data Flow Diagram.''
                     \url{https://www.lucidchart.com/pages/data-flow-diagram}, 2026.
\bibitem{usecase}   Creately. ``Use Case Diagram Tutorial.''
                     \url{https://creately.com/guides/use-case-diagram-tutorial/}, 2026.
\bibitem{gantt}     GanttPRO. ``What Is a Gantt Chart.''
                     \url{https://ganttpro.com/gantt-chart-basics/}, 2026.

\end{thebibliography}

%==============================================================================
%  APPENDICES
%==============================================================================
\appendix

\chapter{Glossary and Supplementary Material}
\label{app:glossary}

\section{Glossary of Terms}
Table~\ref{tab:glossary} defines the retail, tax and engineering vocabulary used in this
report.

\begin{table}[H]
\centering
\caption{Glossary of terms used in this report.}
\label{tab:glossary}
\small
\begin{tabularx}{\textwidth}{l X}
\toprule
\textbf{Term} & \textbf{Meaning} \\
\midrule
POS & Point of sale --- the billing screen and workflow at the till. \\
MRP & Maximum retail price; the VAT-inclusive shelf price. \\
SKU / barcode & Stock-keeping unit code and the scannable EAN used to find items. \\
NBR & National Board of Revenue, Bangladesh's tax authority. \\
Output VAT & The VAT component contained in inclusive prices, back-calculated per slab. \\
Tender & A payment leg of a bill (cash, card, bKash, Nagad or points). \\
GRN & Goods receipt note --- recording delivered stock against a purchase order. \\
Z-report & End-of-shift drawer reconciliation report (day close). \\
Shrinkage & Inventory lost to damage, theft or counting error. \\
bKash / Nagad & Bangladeshi mobile financial services used as tenders. \\
Standing template & A saved weekly restock list cloned into draft purchase orders. \\
PBKDF2 & Password-Based Key Derivation Function~2 (RFC~8018) used for password hashing. \\
Bearer token & The 256-bit session string sent in the Authorization header. \\
RBAC & Role-based access control; permissions derive from the caller's role. \\
GoF & ``Gang of Four'' --- the classic design-patterns catalogue (1994). \\
POJO & Plain old Java object; the domain contains no framework annotations. \\
SPA & Single-page application; one HTML shell with hash-routed views. \\
BigDecimal & Java's exact-decimal arithmetic type used for every money figure. \\
ThreadLocal & Per-thread storage carrying the caller's role for one request. \\
Repository & The persistence interface each of the 11 models is served through. \\
\bottomrule
\end{tabularx}
\end{table}

\section{Key API Endpoints}
\label{app:endpoints}
All endpoints are prefixed with \texttt{/api}; ``Token'' means any authenticated role,
with the minimum role noted. Selected representative endpoints are listed in
Table~\ref{tab:endpoints}.

\begin{table}[H]
\centering
\caption{Key API endpoints.}
\label{tab:endpoints}
\small
\begin{tabularx}{\textwidth}{l l X l}
\toprule
\textbf{Method} & \textbf{Endpoint} & \textbf{Purpose} & \textbf{Access} \\
\midrule
POST   & /auth/login & Authenticate and issue a bearer token & Public \\
GET    & /auth/me & Current caller identity and role & Token \\
POST   & /bill/lines & Add a scanned or picked line to the session bill & Token \\
PUT    & /bill/lines/\{index\} & Edit quantity / weight of a line & Token \\
POST   & /bill/undo & Undo the last bill edit (memento, depth 10) & Token \\
POST   & /bill/tender & Validate and commit the bill (command pipeline) & Token \\
GET    & /products & List / search catalog with view filters & Token \\
POST   & /products & Create a unit / weighed item (factory) & Manager+ \\
DELETE & /products/\{id\} & Delete a catalog item & Admin \\
POST   & /products/\{id\}/adjust & Shrinkage adjustment (damage / loss / theft / count)
& Manager+ \\
GET    & /sales & Sales explorer with filters & Manager+ \\
POST   & /sales/\{receiptNo\}/void & Void a sale with mandatory reason & Manager+ \\
POST   & /sales/\{receiptNo\}/returns & Partial per-line return, pro-rata refund &
Token \\
GET    & /promotions & List promotions & Token \\
POST   & /promotions & Create a promotion & Manager+ \\
POST   & /promotions/validate & Live coupon tester & Manager+ \\
POST   & /customers & Register a loyalty member & Token \\
POST   & /customers/\{id\}/points/adjust & Manual points adjustment & Manager+ \\
POST   & /purchase-orders & Create a draft PO (builder) & Manager+ \\
POST   & /purchase-orders/\{poNo\}/receive & Goods receipt: stock + batch + cost &
Manager+ \\
POST   & /purchase-orders/from-template & Clone a standing template (prototype) &
Manager+ \\
GET    & /reports/\{key\} & Run one of 8 reports, JSON or CSV & varies \\
POST   & /reports/day-close & Close the day and record the Z-report & Manager+ \\
GET    & /audit & Query the activity log & Manager+ \\
GET    & /dev/patterns & Developer Mode pattern catalogue & Developer \\
\bottomrule
\end{tabularx}
\end{table}

\section{Project Timeline}

\begin{figure}[H]
\centering
\includegraphics[width=\textwidth]{timeline.png}
\caption{Project Timeline and Release Milestones.}
\label{fig:timeline}
\end{figure}

\vspace{1.5cm}
\begin{center}
\textit{--- End of Document ---}
\end{center}

\end{document}
